%%%%%%%%%%%%%%%%%%%%%%%%%%%%%%%%%%%%

\section{ポアソン分布}

%%%%%%%%%%%%%%%%%%%%%%%%%%%%%%%%%%%%

\begin{frame}
\frametitle{ポアソン分布}

\begin{itemize}

\item \hl{ポアソン分布}は、固定された母集団において個人が独立である場合、短い単位時間あたりの大きな母集団における稀な事象の数を推定するのに役立つことが多い。

\item ポアソン分布の\hl{レート}とは、主に固定された母集団において単位時間あたりに起こる事象の平均回数のことであり、通常 \mathhl{\lambda} で表される。

\item レートを使って、単位時間内に稀な事象がちょうど $k$ 回観測される確率を記述できる。

\end{itemize}

\vfill

\formula{ポアソン分布}
{
P(稀な事象を $k$ 回観測) = $\frac{\lambda^k e^{-\lambda}}{k!}$、\\
ここで $k$ は0、1、2、……の値をとり、$k!$ は $k$ の階乗を表す。$e \approx 2.718$ は自然対数の底である。\\

この分布の平均と標準偏差はそれぞれ $\lambda$ と $\sqrt{\lambda}$ である。
}

\end{frame}

%%%%%%%%%%%%%%%%%%%%%%%%%%%%%%%%%%%%

\begin{frame}

\dq{ある発展途上国の農村地域では、停電がポアソン分布に従い、週平均2回発生するとする。ある週に停電がちょうど1回だけ発生する確率を計算せよ。}

$\lambda = 2$ が与えられる。

\begin{eqnarray*}
P(\text{週に1回だけ停電}) &=& \frac{2^1 \times e^{-2}}{1!} \\
&=& \frac{2 \times e^{-2}}{1} \\
&=& 0.27
\end{eqnarray*}

\end{frame}

%%%%%%%%%%%%%%%%%%%%%%%%%%%%%%%%%%%%

\begin{frame}

\dq{ある発展途上国の農村地域では、停電がポアソン分布に従い、週平均2回発生するとする。ある特定の\underline{日}に停電が3回発生する確率を計算せよ。}

週単位の停電レートが与えられているが、この問いに答えるには、まずある特定の日の平均停電レートを計算する必要がある：$\lambda_{\text{日}} = \frac{2}{7} = 0.2857$。停電の確率は曜日によらず同じ、すなわち独立であると仮定することに注意する。

\begin{eqnarray*}
P(\text{ある日に3回の停電}) &=& \frac{0.2857^3 \times e^{-0.2857}}{3!} \\
&=& \frac{0.2857^3 \times e^{-0.2857}}{6} \\
&=& 0.0029
\end{eqnarray*}

\end{frame}

%%%%%%%%%%%%%%%%%%%%%%%%%%%%%%%%%%%%

\begin{frame}
\frametitle{ポアソン分布は適用できるか？}

\begin{itemize}

\item 考慮中の事象が稀であり、母集団が大きく、事象が互いに独立に発生する場合、確率変数はポアソン分布に従う可能性がある

\item しかし、事象が実際には独立でない状況も考えられる。例えば、ある夏の間の結婚式の数の確率に関心がある場合、週末の方が結婚式に人気があることを考慮すべきである。

\item この場合、異なる時間帯に異なるレートを持たせることができれば、ポアソンモデルが合理的な場合もある。たとえば、週末のレートを平日より高くモデル化できる。

\item ポアソン分布のレートを別の変数（曜日）に対してモデル化するという考えは、\hl{一般化線形モデル}と呼ばれるより高度な手法の基礎を形成する。本講義の範囲を超えるが、第7章と第8章で線形モデルの基礎を議論する。

\end{itemize}

\end{frame}

%%%%%%%%%%%%%%%%%%%%%%%%%%%%%%%%%%%%

\begin{frame}
\frametitle{練習問題}

\pq{次のどの分布に従う確率変数が正の整数以外の値をとることができるか？}

\begin{enumerate}[(a)]
\item ポアソン分布
\item 負の二項分布
\item 二項分布
\solnMult{正規分布}
\item 幾何分布
\end{enumerate}

\end{frame}
